%!TEX root = ./main.tex

\section[Allocation]{Choosing a Point on the Curve}

% SECTION TIMING: 41:00--52:00 (5 frames).
% SOURCE: "Cross-Domain Dual-Functional OFDM Waveform Design for Accurate
% Sensing/Positioning," IEEE JSAC, 2024. Simulation parameters on the results frame are
% quoted from that work; do not round them.

\begin{frame}{The grid is the budget}
  \vspace{0.3em}
  \begin{columns}[T,onlytextwidth]
    \column{0.50\textwidth}
      \centering
      \begin{tikzpicture}[scale=0.78,transform shape]
        \foreach \x/\y in {0/0,1/0,3/0,4/0,1/1,2/1,5/1,0/2,3/2,4/2,2/3,5/3,0/3,3/3}
          \fill[blue!22,draw=softgray!50] (\x*0.62+0.14,\y*0.52+0.12) rectangle ++(0.58,0.48);
        \foreach \x/\y in {2/0,5/0,0/1,3/1,4/1,1/2,2/2,5/2,1/3,4/3}
          \fill[orange!35,draw=softgray!50] (\x*0.62+0.14,\y*0.52+0.12) rectangle ++(0.58,0.48);
        \draw[-{Latex[length=2mm]},softgray] (0,0) -- (4.05,0)
          node[anchor=north,font=\scriptsize] {time};
        \draw[-{Latex[length=2mm]},softgray] (0,0) -- (0,2.42);
        \node[font=\scriptsize,rotate=90,anchor=south] at (-0.16,1.2) {frequency};
      \end{tikzpicture}

      \vspace{0.4em}
      {\scriptsize\textcolor{ranblue}{$\blacksquare$ communication RE} \quad
       \textcolor{coreorange}{$\blacksquare$ sensing RE}\par}
    \column{0.46\textwidth}
      \small
      \begin{itemize}\itemsep5pt
        \item One RE is one cell.
        \item Each carries symbol and power.
        \item Assign cell, then power.
        \item Two functions, one grid.
      \end{itemize}
  \end{columns}

  \glossline{RE = Resource Element, one subcarrier during one OFDM symbol\\
    PSLR = Peak-to-SideLobe Ratio; RoI = Region of Interest in the delay--Doppler plane}
  \source{https://ieeexplore.ieee.org/document/10683983}{Cross-Domain Dual-Functional OFDM Waveform Design, IEEE JSAC 2024}
\end{frame}

\begin{frame}{Two scores, graded in two rooms}
  \vspace{0.6em}
  \begin{columns}[T,onlytextwidth]
    \column{0.47\textwidth}
      \large\textbf{\textcolor{ranblue}{Communication score}}

      \vspace{0.5em}
      \small
      Achievable data rate, summed over every RE assigned to communication.

      \vspace{0.5em}
      Measured in \textbf{time--frequency}: put bits where the channel is good.
    \column{0.47\textwidth}
      \large\textbf{\textcolor{coreorange}{Sensing score}}

      \vspace{0.5em}
      \small
      PSLR inside the region of interest --- how far the true peak stands above the clutter.

      \vspace{0.5em}
      Measured in \textbf{delay--Doppler}: shape the autocorrelation where targets live.
  \end{columns}

  \vspace{0.8em}
  \takeaway{The same grid is graded twice,\\by two examiners in two different domains.}
\end{frame}

\begin{frame}{Whoever picks first, wins}
  \vspace{0.4em}
  \begin{columns}[T,onlytextwidth]
    \column{0.47\textwidth}
      {\small\textbf{\textcolor{ranblue}{Communication-centric}}\par}
      \vspace{0.4em}
      \small
      \begin{enumerate}\itemsep4pt
        \item Give the good REs to bits.
        \item Then tune leftover powers.
        \item Sensing takes the remainder.
      \end{enumerate}

      \vspace{0.4em}
      {\scriptsize\textcolor{softgray}{Rate stays optimal; PSLR is merely good inside the RoI.}\par}
    \column{0.47\textwidth}
      {\small\textbf{\textcolor{coreorange}{Sensing-centric}}\par}
      \vspace{0.4em}
      \small
      \begin{enumerate}\itemsep4pt
        \item Lock the RoI autocorrelation.
        \item Then tune the rest of the plane.
        \item Bits take the remainder.
      \end{enumerate}

      \vspace{0.4em}
      {\scriptsize\textcolor{softgray}{PSLR is locally perfect; rate approaches, not reaches, optimal.}\par}
  \end{columns}

  \vspace{0.7em}
  \qa{Both designs claim ``almost no loss''. How can that be?}
     {Because the constraint is only enforced inside the \textbf{region of
      interest}. Narrow the room you must be perfect in, and the rest of the grid
      is free for the other function.}
\end{frame}

% VERIFY: parameters quoted from the JSAC 2024 paper's simulation section.
\begin{frame}{What the simulation actually showed}
  \vspace{0.5em}
  \begin{columns}[T,onlytextwidth]
    \column{0.47\textwidth}
      {\small\textbf{\textcolor{ranblue}{Communication-centric}}\par}
      \vspace{0.4em}
      {\small Keeps the optimal achievable rate, and still shows good
       autocorrelation inside the RoI.\par}
    \column{0.47\textwidth}
      {\small\textbf{\textcolor{coreorange}{Sensing-centric}}\par}
      \vspace{0.4em}
      {\small Approaches the optimal rate while guaranteeing a locally perfect
       autocorrelation.\par}
  \end{columns}

  \vspace{0.9em}
  {\scriptsize\textcolor{softgray}{Setup: carrier 240\,GHz, subcarrier spacing 240\,kHz,
   32 consecutive OFDM symbols, 128 subcarriers, symbol length 4.1470\,\textmu s,
   cyclic prefix 1.0368\,\textmu s.}\par}

  \vspace{0.7em}
  \ask{240 GHz and 240 kHz spacing.\\What kind of link is being assumed here?}
\end{frame}

% SOURCE: multiuser ODDM downlink; sum-rate against CRLB, decomposed into two
% subproblems. From the source talk's resource-allocation section.
\begin{frame}{Many users, one echo}
  \vspace{0.4em}
  \begin{columns}[T,onlytextwidth]
    \column{0.45\textwidth}
      \centering
      \begin{tikzpicture}[scale=0.9,transform shape]
        \draw[-{Latex[length=2mm]},softgray] (0,0) -- (3.6,0)
          node[anchor=north east,font=\scriptsize] {CRLB};
        \draw[-{Latex[length=2mm]},softgray] (0,0) -- (0,2.7)
          node[anchor=south,font=\scriptsize] {sum-rate};
        \draw[very thick,draw=softgreen] (0.5,2.25) .. controls (1.4,2.0) .. (3.1,1.45);
        \node[font=\scriptsize,text=softgreen,anchor=south west] at (1.6,2.02) {ODDM};
        \draw[very thick,draw=softgray!70,dashed] (0.5,1.55) .. controls (1.4,1.3) .. (3.1,0.75);
        \node[font=\scriptsize,text=softgray,anchor=north west] at (1.6,1.22) {OTFS};
      \end{tikzpicture}
    \column{0.51\textwidth}
      \small
      \begin{itemize}\itemsep5pt
        \item Power is split over subcarriers.
        \item Other users are interference.
        \item CRLB bounds the estimate.
        \item The problem splits in two.
      \end{itemize}

      \vspace{0.4em}
      {\small\textcolor{softgreen}{\textbf{ODDM sits above OTFS across the whole
       range.}}\par}
  \end{columns}

  \glossline{SINR = Signal to Interference-plus-Noise Ratio \quad
    Fisher information $\rightarrow$ CRLB by inversion}

  \vspace{0.3em}
  \takeaway{Sparser echoes mean less interference,\\and less interference is rate you keep.}
\end{frame}
